On 3 January 2019, the Chinese spacecraft Chang'e-4 landed on the far side of
the Moon in the area of the von K\'arm\'an crater, which was named after the
world famous Hungarian-born physicist Theodore von K\'arm\'an.

As the Earth remains below the horizon of the spacecraft all the time, a relay
station is also necessary for the communication with mission control on the
Earth. For this purpose, The Chinese Space Agency launched a spacecraft,
Queqiao, which was placed into a halo orbit around the outer Lagrangian point
of the Earth--Moon system, $L_2$, the far side of the Moon.

Calculate the distance ($h$) of this satellite above the surface of the Moon.
The Moon's orbit should be considered as a perfect circle with a radius of
$R = 384\,400$\,km. Neglect perturbations from the Sun and other planets.

\emph{Hint:} You can use the following approximation:
$1/(1+x)^2 \approx 1 - 2x$, if $|x| \ll 1$.
